% Generated draft: requires mathematical and source review.
\documentclass[11pt]{article}
\usepackage[margin=25mm]{geometry}
\usepackage{fontspec}
\setmainfont{Latin Modern Roman}
\usepackage{amsmath,amssymb,mathtools,unicode-math}
\setmathfont{Latin Modern Math}
\usepackage{enumitem,needspace,xurl,hyperref}
\setlength{\parindent}{0pt}
\setlength{\parskip}{5pt}
\setlength{\emergencystretch}{3em}
\newcommand{\N}{\mathbb{N}}
\newcommand{\Z}{\mathbb{Z}}
\newcommand{\Q}{\mathbb{Q}}
\newcommand{\R}{\mathbb{R}}
\newcommand{\C}{\mathbb{C}}
\newcommand{\F}{\mathbb{F}}
\newcommand{\A}{\mathbb{A}}
\newcommand{\PP}{\mathbb{P}}
\newcommand{\E}{\mathbb{E}}
\newcommand{\eps}{\varepsilon}
\newcommand{\dd}{\,\mathrm{d}}
\begin{document}
\section*{UnsolvedMath ID: 800463}
\textbf{Source problem code:} catalogue entry
\subsection{Definitions and mathematical statement}
Let \((M_j^n,g_j,x_j)\) be a sequence of pointed Riemannian manifolds, and let \((X,d,x)\) be the pointed limit appearing in the source. Write \(R_j\) for the scalar curvature of \(g_j\), \(\mathrm{d}v_j\) for its Riemannian volume measure, and \(B_1(x_j)\) for the metric ball of radius one. Let \(\mathcal H^k\) denote \(k\)-dimensional Hausdorff measure, and write \(\mathcal H^k\!\restriction_A\) for its restriction to a set \(A\). A countably \(k\)-rectifiable set is, up to an \(\mathcal H^k\)-null set, covered by countably many Lipschitz images of subsets of \(\R^k\).

The supplied Conjecture 7 asserts that, whenever
\[
(M_j^n,g_j,x_j)\longrightarrow(X,d,x),\qquad
\operatorname{Ric}_{g_j}\geq-\lambda,\qquad
\operatorname{Vol}_{g_j}(B_1(x_j))\geq v>0,
\]
the scalar-curvature measures converge as follows:
\[
R_j\,\mathrm{d}v_j\longrightarrow
R_X\,\mathcal H^n
+2\theta_{\cdot}\,\bigl(\mathcal H^{n-1}\!\restriction_{\Sigma}\bigr)
+2\pi(1-r_{\cdot})\,\bigl(\mathcal H^{n-2}\!\restriction_S\bigr).
\]
Here the asserted conclusion requires a function \(R_X\colon X\to\R\) with
\[
R_X\in L^1_{\mathrm{loc}}(X,\mathcal H^n),\qquad
R_X\geq C\quad\text{for some }C\in\R,
\]
a countably \((n-1)\)-rectifiable set \(\Sigma\subseteq X\), and coefficients satisfying \(\theta_y\leq1\). The source further requires
\[
2\pi(1-r_{\cdot})\,\bigl(\mathcal H^{n-2}\!\restriction_S\bigr)
\]
to be an \((n-2)\)-rectifiable measure supported on the top stratum of \(\operatorname{Sing}(X)\). The variable \(y\) in \(\theta_y\), and the dots in the displayed densities, distinguish variable points from the distinguished base point \(x\); they retain the source's coefficients \(\theta_x\) and \(r_x\).

This assertion retains unresolved qualifications from the supplied record. The convergence topology for the spaces, the meaning of measure convergence on the varying spaces, the definitions of \(\theta_y\), \(r_y\), \(S\), and the top singular stratum, and the normalization conventions are not supplied. No lower bound for \(\theta_y\) or range for \(r_y\) is stated. The dimension range, completeness and boundary hypotheses, and precise convention for the Ricci lower bound are also unspecified. These omissions prevent a fully determined mathematical formulation; the displayed identity is a statement-only transcription requiring review.
\subsection{Short English statement}
Naber's supplied energy-identity conjecture predicts that scalar-curvature measures of a noncollapsing sequence with a common Ricci lower bound converge to a sum of a locally integrable bulk term and rectifiable terms of codimensions one and two, with exactly the coefficients displayed above. Several geometric coefficients and convergence conventions are undefined in the record. The registry labels the problem open; this label is supplied metadata, not an independently verified assessment.
\subsection{Sources}
\begin{itemize}
\item Supplied local UnsolvedMath registry record: permanent ID \(800463\), exact source problem code \texttt{catalogue entry}, title ``Spaces with Lower Ricci Curvature --- Conjecture 7 (energy identity).'' The record attributes the conjecture to Naber, with proposed year 2020, and labels it open.
\item Naber, \emph{Conjectures and Open Questions on Spaces with Lower Ricci Curvature} (2020), Conjecture 7; supplied source URL: \url{https://arxiv.org/abs/2010.10031}. The record reports extraction by \texttt{arxiv\_latex} and an access date of 2026-10-16. This access information is supplied provenance and was not independently verified here. The supplied attribution/context is ``energy identity.''
\end{itemize}
\end{document}
